Theorem 3. Let $J$ be a topological 1 -sphere in $\mathbf{R}^{2}$. Then $\mathbf{R}^{2}$ - $J$ is not connected.

Proof. Let $\bar{I}$ be a polyhedral 2-cell containing $J$, such that $J \cap \mathrm{Fr} I$ contains exactly two points $P$ and $R$. (Fill in the details for the construction of such an $\bar{I}$.) Then $J$ is the union of two $\operatorname{arcs} A_{1}$ and $A_{2}$, from $P$ to $R$. Take a broken line $B$, from $S$ to $Q$, in $\bar{I}$, and intersecting $\mathrm{Fr} I$ only at $S$ and $Q$. Let $T$ be the first point of $B$ (in the order from $S$ to $Q$ ) which lies in $J$; let $A_{1}$ be the arc from $P$ to $R$ in $J$ that contains $T$; and let $A_{2}$ be the other arc from $P$ to $R$ in $J$. Let $X$ be the last point of $B$ that lies in $A_{1}$. (See Figure 4.2.)

\begin{figure}[h]
\begin{center}
  \includegraphics[alt={},max width=\textwidth]{c570eb93-f29d-4347-b75d-c251ce2386f6-032_498_738_1692_693}
\captionsetup{labelformat=empty}
\caption{Figure 4.2}
\end{center}
\end{figure}

Lemma 1. $A_{2}$ contains a point of $B$, following $X$ in the order from $S$ to $Q$ on B.

Proof. Suppose not. Let $B_{1}$ be the $\operatorname{arc} S T$ in $B$; let $B_{2}$ be the $\operatorname{arc} T X$ in $A_{1}$; and let $B_{3}$ be the $\operatorname{arc} X Q$ in $B$. Then $B_{1} \cup B_{2} \cup B_{3}$ is an $\operatorname{arc} S Q$, in $\bar{I}-A_{2}$. Therefore $S$ and $Q$ lie in the frontier of the same component of $I-A_{2}$; and this contradicts the preceding theorem.

Now let $Y$ be the first point of $B$ that lies in $A_{2}$ and follows $X$ in $B$ (in the order from $S$ to $Q$ ). Let $Z$ be any point between $X$ and $Y$ in $B$.

Lemma 2. $Z$ lies in a bounded component of $\mathbf{R}^{2}-J$.\\
Proof. Suppose not. Then there is a broken line $B_{1}$, from $Z$ to a point $W$ of Fr $I$, with $B_{1} \subset \mathbf{R}^{2}-J$. We may suppose that $B_{1} \cap \mathrm{Fr} I=W$, since otherwise a shorter broken line would have the same properties. Consider first the case in which $W$ and $S$ lie in the same component of $\operatorname{Fr} I- \{P, R\}$. (See Figure 4.3.)

\begin{figure}[h]
\begin{center}
  \includegraphics[alt={},max width=\textwidth]{c570eb93-f29d-4347-b75d-c251ce2386f6-033_497_742_1018_700}
\captionsetup{labelformat=empty}
\caption{Figure 4.3}
\end{center}
\end{figure}

$B$ contains an arc $B_{2}$, from $Z$ to $Q$, not intersecting $A_{1}$. It follows that $W$ and $Q$ are in the frontier of the same component of $I-A_{1}$, and this contradicts the preceding theorem.

Suppose, second, that $W$ and $Q$ lie in the same component of $\mathrm{Fr} I$ $\{P, R\}$. (See Figure 4.4.) As before, let $T$ be the lowest point of $B=S Q$

\begin{figure}[h]
\begin{center}
  \includegraphics[alt={},max width=\textwidth]{c570eb93-f29d-4347-b75d-c251ce2386f6-033_499_701_1845_715}
\captionsetup{labelformat=empty}
\caption{Figure 4.4}
\end{center}
\end{figure}

that lies on $A_{1}$. Form the union of the $\operatorname{arcs} B_{1}, Z X \subset Q S, X T \subset A_{1}$, and $T S \subset S Q$. It follows that $W$ and $S$ lie in the frontier of the same component of $I-A_{2}$, and this contradicts the preceding theorem, as before.

Evidently Lemma 2 proves Theorem 3. \(\square\)
